\subsubsection{アーキテクチャレビュー}
アーキテクチャレビューは、アーキテクチャの状態、品質、リスクを確認するための有効な方法である。レビューは、非公式な専門家レビューとして行われることもあれば、チェックリストに基づいて構造化されることもある。

より能動的なレビューでは、単に「この項目を確認したか」と聞くのではなく、レビュアが特定の活動を行って情報を得る。たとえば、「障害が発生したときに、どの構成要素がどの順で影響を受けるかを説明する」「利用者数が 10 倍になったときに、どの構成要素がボトルネックになるかを示す」といった活動である。

\par\smallskip\noindent\refstepcounter{table}\label{tab:ch02-06}\textbf{表 \thetable: 第02章 ソフトウェアアーキテクチャ 表06}\par\smallskip
\begin{longtable}{P{0.25\linewidth}P{0.67\linewidth}}
\toprule
レビュー観点 & 確認すること \\
\midrule
要求対応 & ASR に答える構成や判断があるか。 \\
品質特性 & 性能、可用性、変更容易性、安全性、セキュリティなどが説明されているか。 \\
トレードオフ & ある品質を優先したことで、別の品質にどのような影響があるか。 \\
リスク & 未検証の前提、難しい技術、外部依存、単一障害点があるか。 \\
説明可能性 & 重要な意思決定に根拠が残っているか。 \\
実装適合性 & 実装がアーキテクチャに従っているか、または逸脱が管理されているか。 \\
\bottomrule
\end{longtable}
